\section{The unique field and the modular connection of PT}
\label{sec:field_connection}

Before evaluating position C in the PT framework
(Section~\ref{sec:holonomy_primitive}), three structural facts
must be established: that PT admits a unique dynamical degree of
freedom (\S~\ref{ssec:BA0_en}), that this degree of freedom
naturally carries a modular gauge connection
(\S~\ref{ssec:gauge_connection_en}), and that this connection
induces the holonomy angles $\theta_p$ through an explicit
algebraic identity (\S~\ref{ssec:T6_en}). All three facts are
established in the monograph~\cite{PT_Monograph} and reproduced
here without proof.

\subsection{The unique dynamical field BA0/T0}
\label{ssec:BA0_en}

The closure theorem T0 of chapter~25 of the
monograph~\cite{PT_Monograph} establishes that the sequence of
gaps between consecutive primes,
\begin{equation}
  g_n \;:=\; p_{n+1} - p_n \qquad (n \geq 1),
  \label{eq:g_n_def_en}
\end{equation}
is the \emph{unique} dynamical degree of freedom of the
Eratosthenes sieve, under the structural conditions U1--U4
(antidiagonality of $T_3$, GFT identity, canonical geometry of
$\DKL$ and Fisher, fixed point $\mustar = 15$). Postulate BA0 of
the monograph --- that $\{g_n\}_{n \geq 1}$ is \emph{the}
dynamical field of the theory --- is therefore promoted from
postulate to theorem, and the theory admits \emph{zero} dynamical
axioms beyond the standard tools of number theory.

This is what distinguishes PT radically from classical gauge
theories: whereas the latter must postulate a vector field $A_\mu$
as the primitive structure of electrodynamics, PT derives
\emph{all} of the gauge structure --- including the very
existence of a connection --- from the single scalar field
$\{g_n\}$ and the arithmetic coherence constraints. The gauge
connection will be, in what follows, a \emph{derived} object, not
a postulate.

\subsection{The modular gauge connection}
\label{ssec:gauge_connection_en}

For any modulus $m \geq 2$, the additive recurrence
\begin{equation}
  r_{n+1}^{(m)} \;=\; \bigl(r_n^{(m)} + g_n\bigr) \bmod m
  \label{eq:gauge_recursion_en}
\end{equation}
expresses that the trajectory of modular residues
$r_n^{(m)} = p_n \bmod m$ is entirely determined by the initial
datum and the sequence $\{g_n\}$. This recurrence is the
\emph{parallel transport equation} associated with the modular
connection of PT: at each step $n$, the residue is translated by
the ``gauge gain'' $g_n \bmod m$, and one reads the trajectory as
a transport along the ``arithmetic time'' $n$.

The gauge connection theorem~\cite[ch.~1, Thm.~gauge]{PT_Monograph}
formalises this reading: the
recurrence~\eqref{eq:gauge_recursion_en} defines, for every
squarefree integer $m = q_1 \cdots q_k$ that is a product of
distinct primes, a connection on the trivial bundle
$\Omega \times (\mathbb{Z}/m\mathbb{Z}) \to \Omega$ above the
path space $\Omega = \mathbb{N}^{\mathbb{N}}$. By the CRT
decomposition
\begin{equation}
  \mathbb{Z}/m\mathbb{Z}
  \;\cong\;
  \prod_{i=1}^{k} \mathbb{Z}/q_i\mathbb{Z},
  \label{eq:CRT_iso_en}
\end{equation}
this connection splits into $k$ independent components, one per
prime $q_i$, and one can treat each factor
$\mathbb{Z}/q_i\mathbb{Z}$ as a discrete circle of $q_i$ points
carrying its own gauge connection. It is this decomposition that
makes the modular connection of PT structurally simpler than a
classical connection on $\mathbb{R}^4$: it is, by construction,
a tensor product of finite-dimensional objects.

\subsection{The holonomy identity T6}
\label{ssec:T6_en}

The holonomy theorem T6~\cite[ch.~6]{PT_Monograph} gives the
explicit expression of the holonomy angle associated with parallel
transport along a cycle of the modular connection. For each prime
$p$ and each scale $\mu > 0$,
\begin{align}
  \delta_p(\mu) &\;:=\; \frac{1 - q(\mu)^p}{p},
  &
  \cos \theta_p(\mu) &\;:=\; 1 - \delta_p(\mu),
  \label{eq:delta_theta_en}
\\
  \sinp{p}(\mu) &\;=\; \delta_p(\mu)\bigl(2 - \delta_p(\mu)\bigr),
  &
  q(\mu) &\;\in\; \{\qplus,\; q_-\},
  \label{eq:sin2_theta_en}
\end{align}
where $\qplus(\mu) = 1 - 2/\mu$ and $q_-(\mu) = e^{-1/\mu}$ are
the two canonical values of $q$ singled out by the
maximum-entropy bifurcation L0~\cite[ch.~3]{PT_Monograph}. For the
holonomy questions treated here we adopt the coupling branch
$q = \qplus$, which describes the electromagnetic sector and the
leptons.

Identity~\eqref{eq:sin2_theta_en} is purely algebraic --- it uses
only the trigonometric definition
$\sin^2 \theta = 1 - \cos^2 \theta = 1 - (1 - \delta)^2 =
\delta (2 - \delta)$ --- but it has a striking consequence: the
usual trigonometry, and with it the constant $\pi$, \emph{emerges}
from the modular dynamics without having been introduced at the
start. The continuous circle is never postulated; it appears as
the continuous limit of the discrete circle
$\mathbb{Z}/p\mathbb{Z}$ as $p \to \infty$, and $\pi$ appears via
the Riemann zeta function $\zeta(2) = \pi^2/6 =
\prod_p (1 - 1/p^2)^{-1}$, which directly relates the measure of
the continuous circle to the sieve of primes.

\paragraph{Structural summary.}
The three elements established in this section sketch an
architecture in which the gauge connection and the holonomy
angle are \emph{derived} objects, computable from the single
scalar field $\{g_n\}$, and not independent postulates. Position
C of the foundational debate --- holonomy as primitive --- takes,
under PT, a substantially stronger sense than it has in field
theory on manifold: it is not an option of interpretation but
the unique outcome of a structural construction that admits no
other choice. This is the status that the following section
formalises.
